% TOP_PROBLEM: {"id": 3391, "title": "Magic square of squares", "category_id": 1, "category": {"id": 1, "name": "number_theory", "display_name": "Number Theory", "description": "Properties of integers, prime numbers, Diophantine equations.", "slug": "number-theory", "order_index": 1, "created_at": "2026-07-31T15:26:25.670Z"}, "status": "open", "problem_number": "OPG-16570", "set_id": 12}
% FIRST_POSTED: 2026-10-01
% ENTRY_KIND: statement_only
% UnsolvedMath ID 3391; source catalogue code OPG-16570
% !TeX program = lualatex
% Definition and sources only; this file is not an LLM solution attempt.
\documentclass[11pt,a4paper]{article}
\usepackage[margin=25mm]{geometry}
\usepackage{fontspec}
\setmainfont{lmroman10-regular.otf}[BoldFont=lmroman10-bold.otf,
 ItalicFont=lmroman10-italic.otf,BoldItalicFont=lmroman10-bolditalic.otf]
\usepackage{amsmath,amssymb,mathtools}
\usepackage{unicode-math}
\setmathfont{latinmodern-math.otf}
\AtBeginDocument{\let\setminus\smallsetminus}
\usepackage{xurl,hyperref}
\hypersetup{colorlinks=true,urlcolor=blue}
\setcounter{secnumdepth}{0}
\setlength{\parindent}{0pt}
\setlength{\parskip}{5pt}
\setlength{\emergencystretch}{3em}
\begin{document}
\section*{Magic square of squares}\label{3391}
{\small UnsolvedMath ID 3391 · OPG-16570 · Number Theory · OpenGarden}
\subsection{Definitions and mathematical statement}
A $3\times3$ magic square is an array $(a_{ij})_{1\le i,j\le3}$
for which the three row sums, three column sums, and two principal
diagonal sums all equal a common integer $M$. This problem requires
each entry to be a positive integer square and all nine entries to
be pairwise distinct. In explicit variables, ask whether there are
$x_{ij}\in\mathbb Z_{>0}$ with pairwise distinct $x_{ij}^2$ and
\[
 \sum_jx_{ij}^2=M\ (i=1,2,3),\qquad
 \sum_ix_{ij}^2=M\ (j=1,2,3),
\]
\[
 x_{11}^2+x_{22}^2+x_{33}^2
       =x_{13}^2+x_{22}^2+x_{31}^2=M.
\]
The target is existence of one such array, or a nonexistence proof.
The positive-square convention is the standard version intended
here; allowing zero is a separate boundary variant. Distinctness
rules out constant arrays and other trivial repeated-entry examples.

Rotations and reflections of an array preserve the conditions.
Multiplying every entry by the same positive square also preserves
them. Neither symmetry has to be quotiented out for the existence
question. Checking only rows and columns gives a semimagic square
and leaves two required diagonal equalities unverified. Both
diagonals must be checked for a proposed example.
\subsection{Short English statement}
Is there a three-by-three magic square whose nine entries are distinct positive perfect squares?
\subsection{Sources}
\begin{itemize}
\item[S1] ULAM.AI and the UnsolvedMath Contributors, supplied problems.json, record 3391 (OPG-16570), OpenGarden collection. Catalogue identity and source transcription; CC BY 4.0.
\url{https://huggingface.co/datasets/ulamai/UnsolvedMath}.
\item[S2] Open Problem Garden, Magic square of squares. Original problem and discussion; the mathematical formulation below is expanded from the supplied source record.
\url{http://www.openproblemgarden.org/op/magic_square_of_squares}.
\end{itemize}
\end{document}
